\section{Step 74 Separate selector and sample cursors}
Fill \texttt{row\_cursor} at \texttt{filter\_row}. Load the selector into \texttt{\$t0}, advance \texttt{\$a0} past that selector, and set \texttt{\$t1} to the first destination sample. Set \texttt{\$t2} to the corresponding upper sample and \texttt{\$t3} to the exclusive end of this row’s sample bytes. The selector remains in its original buffer position.

On row zero of the worked image, the destination is buffer plus one and the end is buffer plus seven. The upper pointer is the private zero row. These are sample cursors: neither points to the selector in the preceding row. That distinction matters on every row after the first.

The supplied dispatch skips prediction for selector zero. None means its residuals already are the reconstructed samples, so no sample store is required. Other selectors choose the scalar byte loop, except that the MSA entry may select the Up vector loop. Keep the cursor setup common to these paths so each eventually reaches the same row-completion code.

\paragraph{Observe} Use \texttt{s74-1}; stop before \texttt{row\_cursor\_ready}. Inspect the state described above before continuing.

\Needspace{7\baselineskip}
\paragraph{Check} Where does the upper cursor point?
\begin{itemize}
\item[A.] At the corresponding previous sample, or the zero row.
\item[B.] At the previous selector.
\item[C.] At compressed input.
\end{itemize}
